\documentclass[11pt,a4paper]{article}
\usepackage[T1]{fontenc}
\usepackage[utf8]{inputenc}
\usepackage{lmodern}
\usepackage{amsmath,amssymb}
\usepackage[margin=25mm]{geometry}
\usepackage{longtable,booktabs,array,calc}
\usepackage[hidelinks]{hyperref}
\hypersetup{pdftitle={The self-sufficiency contrast: where Newton's axioms leave motion undefined, the theory wi},pdfauthor={navstokgap research working notes}}
\newcommand{\tightlist}{\setlength{\itemsep}{0pt}\setlength{\parskip}{0pt}}
\setlength{\emergencystretch}{3em}
\setcounter{secnumdepth}{0}
\begin{document}
\section{\texorpdfstring{The self-sufficiency contrast: where Newton's
axioms leave motion undefined, the theory with \(h>0\) defines
it}{The self-sufficiency contrast: where Newton's axioms leave motion undefined, the theory with h\textgreater0 defines it}}\label{the-self-sufficiency-contrast-where-newtons-axioms-leave-motion-undefined-the-theory-with-h0-defines-it}

The thesis in STATE says that Newton's mechanics is the large-action
regime of a theory with \(h>0\), and that it needs \(h\) because its own
axioms do not make a complete theory of motion. This note proves the
first piece of evidence for that claim as one theorem with four parts,
for Kepler's force. The classical flow of \(V=-k/r\) is incomplete:
every zero-angular-momentum datum that moves inward, or is bound,
reaches the centre in finite time, with the explicit collision time
\(t_c=\tfrac\pi2\sqrt{mr_0^3/2k}\) for fall from rest, and Newton's
axioms give no continuation there. The quantum Hamiltonian
\(-\tfrac{h^2}{2m}\Delta-k/r\) is self-adjoint on the domain of the
Laplacian and generates a global unitary group for every state and every
\(h>0\), by Kato's theorem, cited with its hypotheses. The same contrast
holds for the quartic between the fixed-\(h\) unitary limit and the
classical-first construction that needs a branch rule. And the classical
trajectory is a derived object of the \(h>0\) theory: for quadratic
forces the centroid of any state obeys Newton's equation exactly at
every \(h\), and for Kepler's force a coherent state follows the orbit
in the large-action regime for finite times away from the collision set.
Section 5a recovers Hepp's hypotheses and supplies the weighted
localization and first-moment bounds needed for Coulomb. Parts (a), (c)
and the quadratic half of (d) were proved in the first version; (b)
rests on a cited theorem. Theorem 1 was refereed by GPT-6.1 Sol (Codex)
in two bounded passes on 2026-10-01, REFINE then ACCEPT on its statement
and proof as printed; the reports are in
\href{../reviews/2026-10-01-theorem1-self-sufficiency-contrast-referee.md}{\texttt{reviews/2026-10-01-theorem1-self-sufficiency-contrast-referee.md}}.
The October 1 assessment did not endorse Proposition 2. Its replacement
in §5b gives the admissible \(h\)-sets under explicit selection
premises. Dirac's finite-potential-energy selection excludes its
critical endpoint; the critical Schur-form and specified cutoff limits
include it. The Klein--Gordon strict-energy construction has an open
endpoint, whereas its radial principal Friedrichs operator includes that
endpoint. A radial branch is not a theorem about the full Klein--Gordon
evolution. These are conditional action thresholds, not a proof that
every completion of Newton's mechanics needs \(h\). Original note by
Claude Fable, 2026-10-01; revision and proof status recorded below.

\textbf{Dated revision, 2026-10-02 (Kepler localization).} The October 1
version left (L) as an assumption and used centroid language without
moment bounds. Section 5 now proves (L) for a fixed smooth bounded
regularization and the stated Gaussian preparations, using centered
weighted Sobolev estimates and Hardy's inequality. It also proves
position and momentum centroid convergence separately. The primary Hepp
statement is identified below; its norm approximation alone is not used
to control a singular multiplier. These additions are written
derivations, not an independent verification or a new referee
acceptance.

\textbf{Dated correction, 2026-10-02 (selection and endpoints).} The
original Proposition 2 used \(h>h_{\min}\), left equality unsettled, and
nevertheless concluded that any \(h\ge h_{\min}\) was admissible. It
also described a chosen Klein--Gordon branch as defining motion without
an added rule. The replacement below withdraws those unqualified
conclusions. It separates the domains and selection criteria that decide
equality; the original October 1 text remains in Git history and its
referee assessment remains linked above.

Notation: \(h\) is the reduced phase constant, \(m>0\) the mass, \(k>0\)
the force constant, \(r=|x|\) in \(\mathbb R^3\); the planar case is a
remark.

\subsection{1. Statement}\label{statement}

\textbf{Theorem 1 (self-sufficiency contrast for Kepler's force).}

\begin{enumerate}
\def\labelenumi{(\alph{enumi})}
\item
  \emph{Classical incompleteness.} Let \(L=x\times p\) be the angular
  momentum. Every initial datum with \(L=0\) whose radial velocity is
  inward or zero, and every datum with \(L=0\) and negative energy,
  reaches \(r=0\) at a finite time. For a body released from rest at
  distance \(r_0\) the collision time is
  \[t_c=F(r_0):=\frac\pi2\sqrt{\frac{m\,r_0^3}{2k}}.\tag{1}\] For bound
  data, \(E<0\) with \(R=-k/E\): \(t_c\le F(R)\) if the radial velocity
  is inward or zero, and \(t_c\le2F(R)\) if it is outward; for inward
  data with \(E\ge0\), \(t_c\le\tfrac23\sqrt{m/2k}\,r_0^{3/2}\). The
  collision set, on the phase space \(\{(x,p):x\neq0\}\),
  \[\mathcal C=\{L=0\ \text{and}\ (\dot r\le0\ \text{or}\ E<0)\},\] is
  nonempty, invariant under the flow while solutions exist, and of
  measure zero. Along every solution in \(\mathcal C\) the force
  magnitude \(k/r^2\) and the momentum \(|p|\sim\sqrt{2mk/r}\) diverge
  as the centre is approached, so there is no continuation as a solution
  of the original vector field in phase space: the flow of Kepler's
  force is not complete. Newton's axioms, the Laws with the force law,
  assign no motion at the centre; any continuation is an added
  prescription (regularization), and the theorem establishes
  incompleteness, not that no classical completion is possible.
\item
  \emph{Quantum completeness.} For every \(h>0\), \(m>0\) and real
  \(k\), the operator
  \[H_h=-\frac{h^2}{2m}\Delta-\frac kr\qquad\text{on }L^2(\mathbb R^3)\tag{2}\]
  is self-adjoint on the domain \(H^2(\mathbb R^3)\) of the Laplacian
  and bounded below, and \(U_h(t)=e^{-itH_h/h}\) is a strongly
  continuous unitary group defined for every \(t\in\mathbb R\) and every
  initial state in \(L^2(\mathbb R^3)\). This is Kato's theorem
  {[}@Kato1951; metadata{]}, whose hypothesis, that the potential is the
  sum of an \(L^2(\mathbb R^3)\) and a bounded function, is verified in
  §3.
\item
  \emph{The quartic contrast.} For \(V=\lambda q^4\) on the line,
  \(\lambda>0\), Theorem 2(b) of the
  \href{rivero-1998-conjecture-central-forces.md}{1998-conjecture note}
  gives, for every \(h>0\), the strong convergence of the time-sliced
  operators to the unitary evolution \(e^{-iTH_h/h}\) along every
  sequence of partitions, with no extra datum; Theorem 1 there shows
  that the unrestricted halved classical-first construction at zero
  resolution has no limit on a two-cell partition with three stationary
  paths, and Theorem 2(a) there exhibits one repair, a branch selection
  chosen by hand. The contrast is between the fixed-\(h\) construction,
  which exists for all data, and that specific classical-first
  prescription, which needs a datum the action does not supply; it says
  nothing against the quartic Newtonian initial-value problem itself,
  which is globally well posed.
\item
  \emph{The classical trajectory as a derived object.} (d1) For
  \(H=p^2/2m+V\) with \(\deg V\le2\), the expectations
  \(\langle q\rangle_t\), \(\langle p\rangle_t\) of every state with
  finite second moments obey Newton's equations exactly for every
  \(h>0\),
  \[\frac{d}{dt}\langle q\rangle_t=\frac{\langle p\rangle_t}m,\qquad
  \frac{d}{dt}\langle p\rangle_t=-V'(\langle q\rangle_t),\tag{3}\] and a
  coherent state stays a Gaussian whose centroid is the classical
  trajectory. (d2) For Kepler's force, let \(z_0=(x_0,p_0)\) lie on a
  classical orbit whose distance from the centre stays at least
  \(\rho_{\min}>0\) on \([0,T]\). Fix \(0<\delta<\rho_{\min}\) and a
  real \(V_\delta\in C_b^\infty\) equal to \(-k/r\) for \(r\ge\delta\).
  For the normalized coherent preparations of §5a, with width
  \(\ell_h=\sqrt{hL_*/P_*}\) and fixed \(L_*,P_*>0\), the regularized
  and Coulomb evolutions both satisfy the Gaussian norm approximation
  (4) with error \(O(\sqrt h)\), uniformly on \([0,T]\). Section 5a
  proves \[\sup_{[0,T]}\|(V-V_\delta)U_h^{(\delta)}(t)\varphi_h\|_2
  =O(h^{3/2})=o(h),\qquad\text{(L)}\] and separately proves that the
  Coulomb position and momentum centroids are \(q_t+O(\sqrt h)\) and
  \(p_t+O(\sqrt h)\). Here ``large action'' means this fixed-scale limit
  \(h\to0\) at fixed orbit. All constants depend on the fixed interval,
  regularization and preparation; no uniformity in \(T\) or \(\delta\)
  is claimed.
\end{enumerate}

\emph{Conclusion.} The theory with \(h>0\) defines motion for all data
and all times, by (b), and yields Newton's trajectories where they
exist, exactly for quadratic forces and, under the fixed-data hypotheses
of §5a, for Kepler's force away from collisions, by (d), while Newton's
axioms leave motion undefined on the collision set, by (a), and the
unrestricted classical-first construction at zero needs a datum the
action does not supply, by (c).

\subsection{2. Proof of (a)}\label{proof-of-a}

Conservation of \(L\) under a central force is Proposition I of Book I;
for \(L=0\) the motion is on a line through the centre and
\(E=\tfrac m2\dot r^2-k/r\) is conserved. Released from rest at \(r_0\),
\(E=-k/r_0\) and
\[\dot r^2=\frac{2k}m\Big(\frac1r-\frac1{r_0}\Big)=\frac{2k}{m}\,\frac{r_0-r}{r\,r_0},\]
with \(\dot r<0\) for \(t>0\) because the force points inward. Hence
\[t_c=\int_0^{r_0}\frac{dr}{|\dot r|}
=\sqrt{\frac{m r_0}{2k}}\int_0^{r_0}\sqrt{\frac{r}{r_0-r}}\,dr.\] With
\(r=r_0\sin^2\theta\), \(dr=2r_0\sin\theta\cos\theta\,d\theta\) and
\(\sqrt{r/(r_0-r)}=\tan\theta\), the integral is
\(2r_0\int_0^{\pi/2}\sin^2\theta\,d\theta=\pi r_0/2\), which gives (1).

For a bound datum with \(L=0\), \(E<0\) and \(R=-k/E\): \(r\le R\)
throughout, \(|\dot r|^2=\tfrac{2k}m(1/r-1/R)\), and the time from
\(r_0\) to the centre along the inward branch is
\(I(r_0)=\sqrt{m/2k}\int_0^{r_0}\sqrt{rR/(R-r)}\,dr\le I(R)=F(R)\). If
the radial velocity is outward, the body first rises to \(R\) in the
time \(F(R)-I(r_0)\) and then falls in the time \(F(R)\), so its
collision time is \(2F(R)-I(r_0)\le2F(R)\), and it exceeds \(F(R)\): one
fall time does not bound it. For inward data with \(E\ge0\),
\(|\dot r|\ge\sqrt{2k/(mr)}\) gives the time bound
\(\int_0^{r_0}\sqrt{mr/2k}\,dr=\tfrac23\sqrt{m/2k}\,r_0^{3/2}\).
Invariance of \(\mathcal C\): \(L\) and \(E\) are conserved; for \(E<0\)
the condition is then automatic, and for \(E\ge0\) the radial velocity
cannot vanish, since \(\tfrac m2\dot r^2=E+k/r>0\), so its inward sign
is preserved. The condition \(L=0\) is that \(p\) be parallel to \(x\),
which for \(x\neq0\) defines a four-dimensional subset of the
six-dimensional phase space, hence of Liouville measure zero; it is
nonempty. As \(r\to0\) along a solution in \(\mathcal C\), the force
magnitude \(k/r^2\) is unbounded and \(\tfrac m2\dot r^2=E+k/r\) gives
\(|p|\sim\sqrt{2mk/r}\), so the solution leaves every compact subset of
phase space in finite time and admits no continuation as a solution of
the vector field. A continuation through the centre (Levi-Civita or
Kustaanheimo--Stiefel regularization) is an added prescription with an
enlarged notion of solution, in the same position as the branch rule of
(c). \(\square\)

\subsection{3. Proof of (b), with Kato's hypothesis
verified}\label{proof-of-b-with-katos-hypothesis-verified}

Kato's theorem {[}@Kato1951; metadata; the statement used is the one
standardly quoted from it, the primary text was not read in this
session{]}: let \(V=V_1+V_2\) be real with \(V_1\in L^2(\mathbb R^3)\)
and \(V_2\in L^\infty(\mathbb R^3)\). Then for every \(a>0\) the
operator \(-a\Delta+V\) is self-adjoint on \(H^2(\mathbb R^3)\) and
bounded below. The mechanism is the Sobolev bound
\(\|u\|_\infty\le\epsilon\|\Delta u\|+C_\epsilon\|u\|\) for
\(u\in H^2(\mathbb R^3)\), which makes \(V_1\) relatively bounded with
respect to \(\Delta\) with relative bound zero, so that the
Kato--Rellich theorem applies for every value of \(a=h^2/2m\).

Verification for \(V=-k/r\): fix a cutoff radius \(R>0\) and put
\(V_1=-\tfrac kr\mathbf 1_{r<R}\), \(V_2=-\tfrac kr\mathbf 1_{r\ge R}\).
Then \(|V_2|\le|k|/R\), and
\[\int_{\mathbb R^3}|V_1|^2d^3x=k^2\int_0^R\frac{4\pi r^2}{r^2}\,dr=4\pi k^2R<\infty.\]
So the hypothesis holds for every real \(k\), and \(H_h\) in (2) is
self-adjoint on \(H^2(\mathbb R^3)\) for every \(h>0\), \(m>0\). Stone's
theorem then gives the unitary group \(U_h(t)\) for all
\(t\in\mathbb R\), strongly continuous, defined on every
\(\psi\in L^2\), with \(U_h(t)\psi\in H^2\) for \(\psi\in H^2\); states
outside the operator domain have the mild unitary evolution, those
inside satisfy the Schrödinger equation. No initial datum and no time is
excluded; the states supported near \(r=0\) evolve like all others.
\(\square\)

\emph{Planar remark.} In \(\mathbb R^2\) the function \(1/r\) is not
square integrable near the origin, and the operator-sense theorem does
not apply; the planar Coulomb Hamiltonian is defined in the form sense
by the KLMN theorem, since \(1/r\in L^p_{\rm loc}(\mathbb R^2)\) for
\(p<2\) is form-bounded relative to \(-\Delta\) with relative bound
zero. The planar Kepler problem used elsewhere in this repository
therefore also has a global unitary dynamics, by that standard route,
which is not reproved here.

\subsection{4. Proof of (c)}\label{proof-of-c}

The statements are Theorems 1, 2(a) and 2(b) of the
\href{rivero-1998-conjecture-central-forces.md}{1998-conjecture note},
accepted on refereeing there (status line of that note). Theorem 2(b)
proves the product formula for
\(H_h=-\tfrac{h^2}{2m}d^2/dq^2+\lambda q^4\) at every \(h>0\): essential
self-adjointness, the strong limit of the symmetric time slicing along
any partition sequence, and the composition \(\|U_h(T)\|=1\). Theorem 1
proves that the halved classical-first functional on the two-cell
partition has, as the resolution tends to zero, the oscillating value
(4) there with no limit, because the three stationary paths interfere;
Theorem 2(a) restores a limit only by restricting to a neighbourhood of
one path, a datum the action does not contain. Hence for the quartic the
\(h>0\) theory exists for all data without extra rules and the
construction at zero does not. \(\square\)

\subsection{5. Proof of (d)}\label{proof-of-d}

\emph{(d1), quadratic forces.} By Theorem A(a) of the
\href{principia-fifth-postulate.md}{fifth-postulate note}, the
Heisenberg equations are Newton's second law as operator identities for
every real \(h\): \(\dot q=p/m\) and \(\dot p=-V'(q)\). Taking
expectations in a state with finite second moments gives
\(\tfrac{d}{dt}\langle q\rangle=\langle p\rangle/m\) and
\(\tfrac{d}{dt}\langle p\rangle=-\langle V'(q)\rangle\); for
\(\deg V\le2\) the function \(V'\) is affine, so
\(\langle V'(q)\rangle=V'(\langle q\rangle)\), which is (3). The
centroid therefore follows the classical trajectory through
\((\langle q\rangle_0,
\langle p\rangle_0)\) for every \(h>0\) and every such state, with no
approximation. For a coherent state the Wigner function is the Gaussian
transported by the affine flow (proof of Theorem 5 of the
\href{zero-branch-reachability.md}{reachability note}), so the state
stays Gaussian with the classical trajectory as its centroid and
covariance \(M_t\Sigma_hM_t^{\sf T}\). The width at time \(t\) depends
on the flow: for free motion, scaling the mass by \(\lambda\) at fixed
\(h\) and widths sends \(M_{12}=t/m\to0\) and the position width to the
initial \(\sigma\), while for Hooke's force at fixed frequency the width
\(s_t^2=\sigma^2\cos^2\omega t+h^2\sin^2\omega t/(4m^2\omega^2\sigma^2)\)
tends to \(\sigma^2\cos^2\omega t\); the centroid transport is exact in
every case. \(\square\)

\subsubsection{5a. Kepler localization and moments
(2026-10-02)}\label{a.-kepler-localization-and-moments-2026-10-02}

\textbf{Target and assumptions.} Fix \(m,k,L_*,P_*>0\), an initial phase
point \(z_0=(q_0,p_0)\), and \(T<\infty\) such that its Kepler orbit
\(z_t=(q_t,p_t)\) exists on \([0,T]\) and
\(\rho_{\min}:=\min_{[0,T]}|q_t|>0\). Fix \(0<\delta<\rho_{\min}\) and a
real \(V_\delta\in C_b^\infty(\mathbb R^3)\) equal to \(-k/r\) on
\(r\ge\delta\). All these choices are independent of \(h\). Put
\(\ell_h=\sqrt{hL_*/P_*}\), choose \(h_0>0\), and consider
\(0<h\le h_0\) with the normalized preparations
\[\varphi_h(x)=(2\pi\ell_h^2)^{-3/4}
\exp\left[-\frac{|x-q_0|^2}{4\ell_h^2}
          +\frac{i p_0\cdot(x-q_0)}h\right].\] Both Hamiltonians have
domain \(H^2(\mathbb R^3)\). The conclusions to prove are (L), its
Duhamel transfer, and convergence of the two first moments. Neither
singular-multiplier control from an \(L^2\) approximation nor
convergence of an unbounded observable is an assumption.

\emph{Primary Hepp statement.} The published paper {[}@Hepp1974;
passage, §2, Theorem 2.1, p.~268; equations (2.15)--(2.20),
pp.~269--270; Theorem 2.3 and condition (R), p.~271{]} states Theorem
2.1 for one degree of freedom: a real potential, \(C^{2+\epsilon}\) near
the classical orbit, \(\epsilon>0\), with
\(\int |V(x)|^2e^{-\eta x^2}dx<\infty\) for some \(0<\eta<\infty\), and
any self-adjoint realization of the scaled Hamiltonian. The limits are
strong limits of Weyl operators, uniformly on compact time intervals
inside the orbit's existence interval. Equation (2.20) gives the
corresponding norm approximation on Gaussian packets. For finitely many
coordinates, Theorem 2.3 replaces the integrability assumption by (R): a
strongly continuous unitary propagator \(U_h(t,s)\) for all real
\(t,s\), with \(U_h(t,s)U_h(s,r)=U_h(t,r)\), \(U_h(t,s)^*=U_h(s,t)\),
\(U_h(t,t)=I\), and the strong derivative
\[\partial_s U_h(t,s)\psi=\frac ih U_h(t,s)H_h(s)\psi\] on the
sufficiently decaying Gaussian packets specified in (2.27), with
\(H_h(s)\psi\) defined as a differential operator. Theorem 2.3 allows
finitely many particles with vector and scalar potentials locally
\(C^{2+\epsilon}\) along the orbit; here the vector potential is zero.

The three-dimensional \(V_\delta\) satisfies (R) because it is bounded
and its Gaussian packets belong to \(H^2\). The approximation is
\[\sup_{0\le t\le T}\|U_h^{(\delta)}(t)\varphi_h
-e^{iS_t/h}\varphi^h_{z_t,M_t}\|_2=O(\sqrt h),\tag{4}\] where
\(S_t=\int_0^t(p_s\cdot\dot q_s-H_{\rm cl}(z_s))ds\) and the Gaussian is
obtained with the continuously lifted quadratic propagator of the
linearized flow \(M_t\). Its translation convention uses
\(p_t\cdot(x-q_t)\) in the phase. To see the rate in this smooth case,
expand \(V_\delta\) through second order at \(q_t\): the cubic remainder
has norm at most \(B_3\ell_h^3\||y|^3g_t\|_2/6\) on the centered
quadratic Gaussian \(g_t\). Duhamel divides this by \(h\) and integrates
over \([0,T]\); the Gaussian moment is bounded on that fixed interval.
This is a direct specialization of Hepp's proof, not an estimate of the
singular multiplier on the unknown norm remainder.

\textbf{Centered Sobolev bound.} Let
\(u_h(t)=U_h^{(\delta)}(t)\varphi_h\) and write exactly
\[u_h(t,x)=\ell_h^{-3/2}e^{i[S_t+p_t\cdot(x-q_t)]/h}a_h(t,y),
\qquad y=(x-q_t)/\ell_h.\] Hamilton's equations cancel the constant and
linear terms, so \[i\partial_ta_h=G_h(t)a_h,\quad
G_h=-a_*\Delta_y+W_h(t,y),\quad a_*=\frac{P_*}{2mL_*},\]
\[W_h=\frac{V_\delta(q_t+\ell_h y)-V_\delta(q_t)
-\ell_h\nabla V_\delta(q_t)\cdot y}{h}.\] Set
\(B_j=\sup_x\|\nabla^jV_\delta(x)\|\) (the Euclidean tensor norm),
\(\ell_0=\sqrt{h_0L_*/P_*}\), and, for integer \(s\ge1\),
\[b_s=\frac{L_*}{P_*}
\max_{2\le j\le\max(2,s)}\ell_0^{\,j-2}B_j.\] Taylor's integral formula
gives \(|\nabla_yW_h|\le b_s|y|\), and
\(|\partial_y^\gamma W_h|\le b_s\) for \(2\le|\gamma|\le s\). All these
constants are uniform for \(0<h\le h_0\) and \(0\le t\le T\).

With \(D_y=-i\nabla_y\), define
\[E_s(a)=\sum_{|\alpha|+|\beta|\le s}\|y^\alpha D_y^\beta a\|_2^2,
\quad N_s=\binom{s+6}{6},\quad
c_s=a_*s(s+1)+b_s(2^s-1+2s).\] The commutator of \(y^\alpha D_y^\beta\)
with the kinetic term has norm bounded by \(a_*s(s+1)\sqrt{E_s}\). For
the potential term use \[[D_y^\beta,W_h]=
\sum_{0<\gamma\le\beta}\binom{\beta}{\gamma}
(-i)^{|\gamma|}(\partial^\gamma W_h)D_y^{\beta-\gamma}.\] When
\(|\gamma|=1\), the bound \(b_s|y|\) increases the position degree by
one and decreases the derivative degree by one. Bounding \(|y|\) by
\(\sum_j|y_j|\) gives \(3s b_s\sqrt{E_s}\). The remaining binomial
coefficients sum to at most \(2^s-1-s\). No term has degree larger than
\(s\). Thus \[\frac{d}{dt}E_s(a_h)\le2c_s\sqrt{N_s}\,E_s(a_h),\qquad
E_s(a_h(t))\le E_s(a_0)e^{2c_s\sqrt{N_s}T},\tag{4a}\] where
\(a_0=(2\pi)^{-3/4}e^{-|y|^2/4}\). For fixed \(h\) the calculation is
justified on Schwartz packets: the integral equation for a bounded
smooth potential, with position and derivative cutoffs, gives the same
finite-seminorm energy estimates and propagation of Schwartz regularity;
letting the cutoffs tend to infinity proves (4a). Equivalently, one
first obtains these estimates for smooth truncated data and passes by
the \(L^2\) evolution and lower semicontinuity. The centered estimates,
unlike the preliminary fixed-\(h\) ones, have the uniform constants
displayed above.

For use below define
\[A_s(T)=3^{s/2}\sqrt{E_s(a_0)}e^{c_s\sqrt{N_s}T}.\] Expanding
\(|y|^{2n}=(y_1^2+y_2^2+y_3^2)^n\) shows that both \(\||y|^na_h(t)\|_2\)
and \(\||y|^n\nabla_ya_h(t)\|_2\) are at most \(A_{n+1}(T)\).

\textbf{Proof of (L).} Put \(r_c=(\delta+\rho_{\min})/2\) and
\(d=\rho_{\min}-r_c>0\). Choose a smooth radial cutoff \(\chi\) equal to
one for \(r\le\delta\), zero for \(r\ge r_c\), with \(0\le\chi\le1\) and
\(\|\nabla\chi\|_\infty\le2/(r_c-\delta)=2/d\). On its support
\(|x-q_t|\ge d\), hence \(|y|\ge d/\ell_h\). The centered bounds imply,
for every integer \(n\ge1\),
\[\|\chi u_h(t)\|_2\le A_{n+1}(T)(\ell_h/d)^n,\qquad
\|\chi(\nabla-ip_t/h)u_h(t)\|_2
\le A_{n+1}(T)\ell_h^{n-1}/d^n.\] Apply the three-dimensional Hardy
inequality \(\|f/r\|_2\le2\|\nabla f\|_2\) to
\(f=e^{-ip_t\cdot x/h}\chi u_h(t)\); the phase removes the potentially
large carrier momentum from the estimate. Since \(V-V_\delta\) vanishes
outside \(r<\delta\), this gives
\[\sup_{0\le t\le T}\|(V-V_\delta)u_h(t)\|_2
\le\frac{A_{n+1}(T)}{d^n}
\left[(4k/d+B_0)\ell_h^n+2k\ell_h^{n-1}\right].\tag{4b}\] For \(n=4\)
the right side is \(O(h^{3/2})=o(h)\), proving (L). Choosing larger
\(n\) gives any prescribed power of \(h\). The dependence is explicit
through \(T\), \(m\), \(L_*/P_*\), the preparation seminorms,
\(B_j(V_\delta)\), and the separation \(d=(\rho_{\min}-\delta)/2\). No
uniformity as \(T\to\infty\), \(\delta\to0\), or
\(\rho_{\min}-\delta\to0\) is claimed. A bounded family of centered
Schwartz preparations satisfies the same estimate with its uniform
\(E_s(a_0)\) in place of the Gaussian value. For the usual scaled
regularization \(V_\delta(x)=-(k/\delta)v(x/\delta)\), with fixed smooth
bounded \(v\) equal to \(1/|x|\) for \(|x|\ge1\), the derivative
dependence is \(B_j\le k\delta^{-j-1}\|\nabla^jv\|_\infty\). A radius
\(\delta\) alone cannot bound these constants for arbitrary
regularizations.

On the common domain \(H^2\), Duhamel now gives
\[U_h(t)\varphi_h-U_h^{(\delta)}(t)\varphi_h
=-\frac ih\int_0^tU_h(t-s)(V-V_\delta)u_h(s)\,ds.\] Consequently the
supremum of this norm is \(O(h^{(n-3)/2})\) for every \(n\ge4\). In
particular the two evolutions differ by \(O(h^M)\) for every fixed
\(M>0\), with an \(M\)-dependent constant. Equation (4) therefore
transfers to the Coulomb evolution.

\textbf{First moments.} Write \(\psi_h=U_h\varphi_h\) and
\(p=-ih\nabla\). Its conserved form energy \(e_h\) is bounded
independently of \(h\): the initial kinetic energy is
\(|p_0|^2/(2m)+3hP_*/(8mL_*)\), and the elementary spherical Gaussian
integral gives \(\langle\varphi_h,r^{-1}\varphi_h\rangle
=\operatorname{erf}(|q_0|/(\sqrt2\ell_h))/|q_0|\le1/|q_0|\). Let
\(e_*=\sup_{0<h\le h_0}|e_h|<\infty\). Hardy and Young's inequality
give, with \(X_h(t)=\|p\psi_h(t)\|_2\),
\[e_h\ge\frac{X_h(t)^2}{2m}-\frac{2k}{h}X_h(t)
\ge\frac{X_h(t)^2}{4m}-\frac{4mk^2}{h^2},\qquad
X_h(t)\le2\sqrt{m e_*}+\frac{4mk}{h}.\tag{4c}\] The continuity equation,
first with bounded position cutoffs and then by passage to the limit,
gives \[\|x\psi_h(t)\|_2
\le\|x\varphi_h\|_2+\frac1m\int_0^tX_h(s)\,ds=O(1+h^{-1}).\tag{4d}\]
This also proves preservation of the position-moment domain on the fixed
interval. Energy conservation preserves the Coulomb form domain \(H^1\),
so the momentum norm is finite as well. For \(u_h\) the centered bounds
give \(\|(x-q_t)u_h(t)\|_2\le\ell_h A_1(T)\) and
\(\|(p-p_t)u_h(t)\|_2\le(h/\ell_h)A_1(T)\).

For each self-adjoint coordinate or momentum component \(B\) and the two
normalized states, symmetry on its domain yields
\[|\langle\psi_h,B\psi_h\rangle-\langle u_h,Bu_h\rangle|
\le\|\psi_h-u_h\|_2(\|B\psi_h\|_2+\|Bu_h\|_2).\] Use (4b) with \(n=6\):
the norm difference is \(O(h^{3/2})\), while (4c)--(4d) bound the second
factor by \(O(h^{-1})\). Thus, uniformly on \([0,T]\),
\[\langle x\rangle_{\psi_h(t)}=q_t+O(\sqrt h),\qquad
\langle p\rangle_{\psi_h(t)}=p_t+O(\sqrt h).\tag{4e}\] These
first-moment statements use derivatives of \(V_\delta\) through order
seven via (4a). They establish the centroid assertion without assuming
moment convergence from a norm approximation. They give neither
long-time control nor convergence of higher moments. \(\square\)

\subsection{5b. Proposition 2: action thresholds under specified
selections}\label{b.-proposition-2-action-thresholds-under-specified-selections}

Theorem 1(b) gives no positive lower threshold: the non-relativistic
Coulomb evolution exists for every \(h>0\). At finite \(c\) a threshold
depends on which relativistic model, domain and selection premise are
required. Fix \(m,k,c>0\) and \(\nu=k/(hc)>0\). All statements concern
an ideal point Coulomb source. Rescale length by \(\lambda_h=h/(mc)\)
and energy by \(mc^2\); in the formulas for domains below, \(x\) and
\(r\) denote these dimensionless coordinates. This unitary rescaling
leaves \(\nu\) unchanged.

\textbf{Primary sources and their scope.} Schmincke's original example,
p.~348, applies his Theorems 2--3 to the strict range \(|\nu|<1\) and
identifies the selected domain by finite potential energy
{[}@Schmincke1972; passage, pp.~345--348, GDZ text{]}. Wüst's theorem
and corollary require a real \(C^1\) potential dominated by \(\mu/r\)
with \(\mu<1\), together with the specified bounded cutoff and
convergence conditions; the selected domain is
\(D(D_{\min}^*)\cap D(r^{-1/2})\) {[}@Wuest1975; passage, pp.~93--95,
GDZ text{]}. Those results do not settle \(\nu=1\). Nenciu constructs
the finite-kinetic-form-domain extension for a strict subcritical
Coulomb bound {[}@Nenciu1976; passage, introduction and Theorem 5.1,
pp.~246--247, author-posted text{]}; its poorly extracted formulas are
not used here to infer an endpoint.

The precise subcritical domain statement used below is Theorem 1.1 of
Esteban--Lewin--Séré: for real \(V=V_1+V_2+V_3\) with \(|V_1|\le\nu/r\),
\(0\le\nu<1\), \(V_2\in L^3\) and \(V_3\in L^\infty\), there is a unique
extension with \(H^1\subset D(D_V)\subset H^{1/2}\), also uniquely
characterized by \(\int|\Psi|^2/r<\infty\) for every domain vector
{[}@EstebanLewinSere2019; passage, arXiv v7, Theorem 1.1, pp.~9--10{]}.
Their §1.2 and Appendix A include essential self-adjointness at
\(\nu=\sqrt3/2\). Their Theorem 1.11, building on Esteban--Loss
{[}@EstebanLoss2007; passage, Definition 3 and Theorem 4, pp.~3--4{]},
supplies a different critical domain and norm-resolvent cutoff limit.

For Klein--Gordon, Case's bibliographic identity is verified as
\emph{Physical Review} \textbf{80}, 797--806 (1950), DOI
\href{https://doi.org/10.1103/PhysRev.80.797}{10.1103/PhysRev.80.797}
{[}@Case1950; metadata/abstract{]}. Its original body was not accessible
in this revision, so it is not a proof input. We derive the radial
indices below directly and use separate read sources for the domains and
the evolution. Domain and spectral facts for the inverse-square
principal operator come from Dereziński--Richard
{[}@DerezinskiRichard2017; passage, §§2.2--2.3, equations (2.3)--(2.6),
Corollary 2.4 and Theorem 5.5{]}. The actual global Klein--Gordon
construction used here is Langer--Najman--Tretter
{[}@LangerNajmanTretter2008; passage, Example 8.3 and Theorems 8.5--8.6,
pp.~744--747{]}. These are distinct results.

\textbf{Proposition 2 (selection-dependent admissible sets).} The
following sets, including their endpoints, correspond to the premises in
the first column. An admissible set in one row is not an exclusion
theorem for other rows or other completions.

\begin{longtable}[]{@{}
  >{\raggedright\arraybackslash}p{(\linewidth - 4\tabcolsep) * \real{0.3333}}
  >{\raggedright\arraybackslash}p{(\linewidth - 4\tabcolsep) * \real{0.3333}}
  >{\raggedright\arraybackslash}p{(\linewidth - 4\tabcolsep) * \real{0.3333}}@{}}
\toprule\noalign{}
\begin{minipage}[b]{\linewidth}\raggedright
Selection premise
\end{minipage} & \begin{minipage}[b]{\linewidth}\raggedright
Admissible \(h\)-set
\end{minipage} & \begin{minipage}[b]{\linewidth}\raggedright
Result supplied
\end{minipage} \\
\midrule\noalign{}
\endhead
\bottomrule\noalign{}
\endlastfoot
Dirac minimal operator on \(C_c^\infty(\mathbb R^3\setminus\{0\})^4\):
require essential self-adjointness & \([2k/(\sqrt3c),\infty)\) & Unique
closure, global unitary evolution \\
Dirac: require every operator-domain vector to have finite separate
Coulomb potential energy & \((k/c,\infty)\) & Unique distinguished
extension, global unitary evolution \\
Dirac: use that extension below critical coupling and the critical
Schur-domain or specified cutoff-limit prescription at equality &
\([k/c,\infty)\) & Global unitary evolution, including \(h=k/c\) \\
Dirac: permit additional self-adjoint boundary conditions in singular
channels & \((0,\infty)\) & Extensions exist; uniqueness is lost where
the minimal operator is not essentially self-adjoint \\
Klein--Gordon: require strict coercivity of the conserved energy in the
usual \(H^1\oplus L^2\) topology, and use the charge-space generator
below & \((2k/c,\infty)\) & A global charge-space evolution is
constructed; the strict criterion excludes equality \\
Klein--Gordon \(s\)-wave principal radial operator: choose its
Friedrichs extension & \([2k/c,\infty)\) & Principal radial operator and
branch, including the no-log branch at equality; full critical evolution
is not concluded \\
\end{longtable}

\emph{Dirac proof and domains.} In the scaled coordinates the expression
is \(D_\nu=-i\alpha\cdot\nabla+\beta-\nu/r\). In a channel
\(\kappa\in\mathbb Z\setminus\{0\}\) its reduced two-component radial
operator acts in \(L^2((0,\infty),dr)^2\) as
\[d_{\nu,\kappa}=\begin{pmatrix}
1-\nu/r&-\partial_r+\kappa/r\\
\partial_r+\kappa/r&-1-\nu/r
\end{pmatrix}.\] Its origin modes have exponents \(\pm s_\kappa\),
\(s_\kappa=\sqrt{\kappa^2-\nu^2}\). For real \(s_\kappa>0\), the less
regular mode is locally square integrable precisely when
\(s_\kappa<1/2\); for imaginary \(s_\kappa\) both modes are locally
square integrable. At \(s_\kappa=1/2\) its squared norm contains
\(\int dr/r\) and diverges. Infinity is limit point. Taking
\(|\kappa|=1\) gives the inclusive essential self-adjointness threshold
\(\nu\le\sqrt3/2\). The closure has domain \(H^1\) for strict
inequality; at equality its domain is larger than \(H^1\) (Appendix A,
Proposition A.1 of the cited source).

For \(\sqrt3/2<\nu<1\), both modes are in radial \(L^2\). A full spinor
contains the additional factor \(1/r\), so its potential integral is
\(\int_0^\epsilon(|u|^2+|v|^2)dr/r\). Finite potential energy excludes
the \(r^{-s_1}\) coefficient and retains \(r^{s_1}\), giving the unique
domain \(D(D_{\min}^*)\cap D(r^{-1/2})\). At \(\nu=1\) the modes are
constant and logarithmic. Both have divergent separate potential energy.
An extension domain must contain a nonzero boundary mode: setting both
coefficients to zero leaves the minimal, non-self-adjoint domain. Thus
no self-adjoint extension at equality has all domain vectors of finite
separate potential energy. For \(\nu>1\) the lowest channel has
imaginary \(s_1\) and the same integral diverges for every nonzero
boundary mode. This proves the open set in the second row.

The critical prescription instead uses the closed Schur form of
Esteban--Loss. For \(0<\gamma<1\), define on
\(C_c^\infty(\mathbb R^3,\mathbb C^2)\)
\[b_\gamma[\phi]=\int_{\mathbb R^3}\left[
\frac{r}{1+\gamma r}|\sigma\cdot\nabla\phi|^2
+(2-\gamma-r^{-1})|\phi|^2\right]dx.\] Their critical Hardy--Dirac
inequality (7) verifies Assumption (A) with \(c(V)=0\) and
\(\Gamma=\sup V=0\). Let \(\mathcal H_+\) be the form domain of its
closure, equivalently the completion in
\((b_\gamma[\phi]+\|\phi\|_2^2)^{1/2}\). Proposition 2 makes this domain
independent of \(\gamma\) in the stated range. The closure defines the
combined form; its two integrals need not exist separately for a
completed vector. Definition 3 and Theorem 4 give the operator domain
\[\mathcal D_C=\{\Psi=(\phi,\chi)\in L^2(\mathbb R^3)^4:
\phi\in\mathcal H_+,\ D_1\Psi\in L^2\text{ in distributions}\}.\] They
prove self-adjointness and uniqueness under that form-domain premise.
The angular-channel formulation in Esteban--Lewin--Séré, equations
(1.28)--(1.31), admits \(1/r\) spinor singularities and excludes
\((\log r)/r\) in both lowest channels. Its combined operator energy is
finite; its separate potential and kinetic form energies need not be.
Their Theorem 1.11, for \(V\ge-1/r\) and \(\sup V<1\), reports
norm-resolvent convergence for \(V_\epsilon=\max(V,-1/\epsilon)\) or
\(V_\epsilon=(1-\epsilon)V\), as \(\epsilon\downarrow0\). This
establishes the included endpoint under this changed selection premise.

\textbf{Source-formula clarification, 2026-10-02.} The literal compact
weighted-gradient formula printed in (1.32) of arXiv v7 is not used as a
domain definition here: it conflicts with the angular form description
just cited. For a nonzero constant spinor \(\xi\), put
\[\phi(x)=\frac{e^{-r}}r\sigma\cdot\omega_x\xi,
\qquad \omega_x=x/r.\] Its critical radial profile
\(\phi_1=e^{-r}\xi/r\) satisfies \(r\phi_1'+\phi_1=-e^{-r}\xi\), so the
angular form norm in (1.29) is finite. Logarithmic cutoffs of
\(r\phi_1\) have cost \(\int r|\eta_\epsilon'|^2dr=1/\log(1/\epsilon)\)
on \([\epsilon^2,\epsilon]\) and place it in the Schur-form closure. But
\[\sigma\cdot\nabla(r\phi)=e^{-r}(2/r-1)\xi,\] whose squared norm with
weight \(1/[r(1+r)]\) diverges like \(\int_0 dr/r\). Thus that compact
rewrite cannot be taken literally as equivalent to the angular norm. The
source TeX and HTML confirm that this is not a PDF extraction change. We
use the original Esteban--Loss form closure and the explicit angular
conditions; this calculation checks the problematic rewrite, not the
full borrowed theorems independently.

For arbitrary \(\nu>0\), conjugation of the real radial system makes its
two deficiency indices equal. Every limit-circle channel therefore
admits self-adjoint boundary conditions; only finitely many channels
require them for a fixed \(\nu\). Their direct sum gives a global
self-adjoint Dirac realization. Above critical coupling an origin phase
is additional data. Self-adjointness alone consequently gives no
positive lower bound on \(h\).

\emph{Klein--Gordon construction and domains.} In the scaled variables
the equation is \([(i\partial_t-V)^2+\Delta-1]\psi=0\), \(V=-\nu/r\).
Hardy's inequality gives \(\|Vf\|_2\le2\nu\|\nabla f\|_2\) for
\(f\in H^1\). Langer--Najman--Tretter's Theorem 8.5 assumes
\(H^1\subset D(V)\) and a decomposition \(V=V_0+V_1\) with
\[\|V_0f\|_2\le a\|f\|_2+b\|(-\Delta+1)^{1/2}f\|_2,
\quad a+b<1,\] and \(V_1\) relatively compact. Pure Coulomb satisfies it
with \(a=0\), \(b=2\nu<1\), \(V_1=0\). The theorem constructs the
generator \[A\binom{x}{y}=\binom{Vx+y}{(-\Delta+1)x+Vy}\] in the charge
space \(\mathcal K=H^{1/2}\oplus H^{-1/2}\), with domain
\[D(A)=\{(x,y)\in\mathcal K:x\in H^1,\ y\in L^2,
\ Vx+y\in H^{1/2},\ (-\Delta+1)x+Vy\in H^{-1/2}\}.\] The charge pairing
is \([X,X']=\langle x,y'\rangle+\langle y,x'\rangle\), using Sobolev
duality. The group is strongly continuous and unitary for this
indefinite pairing. It gives global mild evolution on \(\mathcal K\) and
classical evolution on \(D(A)\). It is not a probability-unitary group
on scalar \(L^2\). Theorem 8.6 gives the corresponding strong
Klein--Gordon solutions for \(\psi_0\in H^1\), \(\psi_1\in H^{1/2}\) and
\((-\Delta-V^2)\psi_0\in H^{-1/2}\). For the sign convention of the
equation above use \(i\partial_tX=AX\), \(y=(i\partial_t-V)\psi\), and
\(e^{-itA}\); the source's group \(e^{itA}\) also gives this evolution
after reversing the time parameter.

In physical units the conserved energy has spatial part
\[Q_\nu[f]=h^2c^2\|\nabla f\|_2^2+m^2c^4\|f\|_2^2
-k^2\|f/r\|_2^2
\ge(1-4\nu^2)h^2c^2\|\nabla f\|_2^2+m^2c^4\|f\|_2^2.\] Sharpness of
Hardy's constant, with dilations of its near-extremizers to remove the
mass term relative to the gradient, shows that strict coercivity in the
usual \(H^1\) norm holds exactly for \(\nu<1/2\). At equality the
combined form is nonnegative but loses that coercivity; above it the
spatial form is unbounded below. This settles the fifth row's criterion,
not the existence of all possible Klein--Gordon completions.

For a stationary \(s\)-wave \(\psi=e^{-iEt/h}u(r)/r\), direct
substitution gives the pencil
\[-h^2c^2u''+(m^2c^4-E^2)u-\frac{2Ek}{r}u-\frac{k^2}{r^2}u=0.\] The
indices are \(1/2\pm s\), \(s=\sqrt{1/4-\nu^2}\). The principal operator
\(B_\nu=-\partial_r^2-\nu^2/r^2\) starts on \(C_c^\infty(0,\infty)\) in
\(L^2(dr)\), with maximal domain consisting of \(u,u'\) locally
absolutely continuous and \(B_\nu u\in L^2\). For \(0<\nu<1/2\) its
Friedrichs domain sets the coefficient of \(r^{1/2-s}\) to zero. Both
branches have finite \(u(0)\) and are locally \(L^2\); those properties
alone do not select it. At \(\nu=1/2\) the modes are \(\sqrt r\) and
\(\sqrt r\log r\). The Friedrichs domain sets the logarithmic
coefficient to zero, admitting the endpoint for this principal operator.
It permits a nonzero \(\sqrt r\) coefficient whose separate kinetic
integral diverges; the critical combined form is finite: writing
\(u=\sqrt r\,v\) on \(C_c^\infty(0,\infty)\) gives
\(\int(|u'|^2-|u|^2/(4r^2))dr=\int r|v'|^2dr\). A logarithmic cutoff
near zero approximates a constant \(v\) with vanishing form cost. For
\(\nu>1/2\), \(s=i\beta\), self-adjoint extensions require an origin
phase and all have a negative spectrum unbounded below
(Dereziński--Richard, Theorem 5.5(ii)). No semibounded Friedrichs
extension exists there. The lower-order \(1/r\) term of the stationary
pencil does not change the indicial classification.

\textbf{Uncontrolled critical dynamical term.} The no-log radial
selection does not prove a global Klein--Gordon group at \(\nu=1/2\).
The critical combined form does not bound \(\|f/r\|_2\): a cutoff of
\(\sqrt r\) in the radial form domain already has infinite separate
kinetic and inverse-square integrals. Thus the multiplier bound with
relative constant strictly below one fails. On the enlarged critical
form space, the mixed dynamical term \(V\partial_t\psi\), or
equivalently the required control of \(VQ_{1/2}^{-1/2}\), remains
unbounded by the estimates used here. A compatible phase space and
generator domain are still needed for a critical full evolution theorem.
No such theorem is asserted by the sixth row. \(\square\)

\textbf{Physical selection premises.} Dirac self-adjointness ensures
probability conservation, independently of the action-threshold
argument. It does not select among its extensions. Requiring each
separate energy to be finite is an additional regularity premise; finite
total energy does not entail it, as the critical domain shows. The
specified bounded-source and subcritical-coupling limits give
independent mathematical support for the critical Dirac prescription
once those preparations are chosen. They do not prove that all source
regularizations have that limit or that the ideal point source is
mandatory. Klein--Gordon energy stability supports semibounded principal
forms; Friedrichs closure and a particular topology are additional
choices, and charge conservation is not positive probability
conservation. A finite-size source introduces further data and is
outside these point-source thresholds.

The unit \(k/c\) agrees with the
\href{necessity-unit-and-indeterminacy.md}{necessity-unit note} and the
\href{relativistic-kepler-threshold.md}{relativistic Kepler threshold}.
That dimensional agreement supplies neither a selection principle nor
the measured value of \(h\). In particular there is no single
unqualified set \([h_{\min},\infty)\) shared by these models.

\subsection{6. What the theorem establishes and what it does
not}\label{what-the-theorem-establishes-and-what-it-does-not}

Theorem 1 compares an incomplete classical point-force flow with a
global non-relativistic Coulomb unitary evolution. Section 5a proves the
fixed-interval Kepler localization and both first-moment limits under
its stated regularization and coherent-preparation hypotheses.
Proposition 2 assigns endpoints to explicit selection criteria; it does
not derive those criteria from Newton's axioms.

Still open: the planar operator beyond the cited form construction;
critical full Klein--Gordon dynamics under the no-log prescription;
long-time or shrinking-regularization estimates; and the necessity claim
that every completion of Newton's mechanics requires positive \(h\).
Classical regularizations remain possible added prescriptions, and Dirac
self-adjoint extensions exist at every \(h>0\) if boundary data are
allowed. Excluding those routes requires an independently justified
premise. The quartic Newtonian initial-value problem is globally well
posed; its contrast concerns the particular classical-first time-slicing
construction. These limitations are part of the statement, not
consequences to be inferred from a threshold formula.

\subsection{7. Consequence for STATE}\label{consequence-for-state}

The fixed-\(T\), fixed-\(\delta\) localization obligation (L) and the
Kepler centroid obligation are closed by the written derivation in §5a.
Replace the former unqualified action interval by Proposition 2's
selection-dependent sets: Dirac finite-potential selection
\((k/c,\infty)\), Dirac critical-form/cutoff selection \([k/c,\infty)\),
and Klein--Gordon strict-energy construction \((2k/c,\infty)\); its
radial principal Friedrichs endpoint is included, with full critical
dynamics still open. Next: decide which selection premises can be
required independently of the desired action gap, then address the
exclusion of completions without \(h\). The October 1 referee acceptance
applies to the original Theorem 1 as then printed. The October 2
localization, moment proof and replacement Proposition 2 are written,
source-supported additions by Codex and have not received independent
verification or referee acceptance.
\end{document}
